\documentclass[../../../include/open-logic-section]{subfiles}

\begin{document}

\olfileid{his}{set}{mythology}

\olsection{أسطورة أكثر منها تاريخًا؟}

حين ننظر إلى هذه الأحداث بعد أكثر من قرن، يجب أن نحذر من التسليم بهذه
الأحكام من دون تمحيص. فقد كانت النتائج جديدة ومثيرة ومدهشة بالتأكيد.
لكن إلى أي حد كانت صادمة حقًا؟ وهل بيّنت فعلًا أنه ينبغي لنا ألا نعتمد
على الحدس الهندسي؟

في مسألة الصدمة، يشير~\citet{Gouvea2011} إلى أن ملاحظة كانتور الشهيرة
إلى ديدكند، «\emph{je le vois, mais je ne le crois pas}»، قد اقتُطعت
إلى حد بعيد من سياقها. وفيما يأتي جزء أكبر من ذلك السياق (نقلًا عن
\citeauthor{Gouvea2011}):
\begin{quote}
أرجو أن تعذر حماستي لهذا الموضوع إذا أثقلت عليك بكثرة ما أطلبه من
لطفك وصبرك؛ فالرسائل التي بعثت بها إليك مؤخرًا غير متوقعة وجديدة حتى
بالنسبة إليّ، إلى حد أنني لا أستطيع الاطمئنان حتى أحصل منك، يا صديقي
المكرم، على حكم في صحتها. وما دمت لم توافقني، فلا يسعني إلا أن أقول:
\emph{je le vois, mais je ne le crois pas.}
\end{quote}
كان كانتور يعلم أن نتيجته «غير متوقعة وجديدة إلى هذا الحد». لكن من
المشكوك فيه أنه عدها يومًا \emph{غير قابلة للتصديق}. فكما يشير
\citeauthor{Gouvea2011}، لم يكن يفعل إلا أن يطلب من ديدكند التحقق من
البرهان الذي قدمه.

أما مسألة الحدس الهندسي، فقد نشر بيانو منحناه المالئ للفضاء من دون أي
رسوم. لكن حين نشر هيلبرت منحناه شرح غرضه: سيقدم للقراء طريقة واضحة
لفهم نتيجة بيانو إذا «استعانوا بالحدس الهندسي الآتي»؛ ثم أورد سلسلة
من \emph{الرسوم} شبيهة تمامًا بالرسوم الواردة في
\olref[his][set][pathology]{sec}.

وعلى نحو أعم، مع أن الرسوم لم تعد شائعة كثيرًا في البراهين المنشورة،
فلا سبيل إلى تجاهل أن الرياضيين يستخدمونها \emph{كثيرًا} عند إثبات
الأشياء. (وبعبارة تقريبية: يعرف الرياضيون الماهرون متى يمكنهم الاعتماد
على الحدس الهندسي.)

وخلاصة القول: لا تصدق التهويل، أو على الأقل لا تسلّم به من دون تمحيص.
ولمزيد من التفصيل، يمكنك قراءة~\citet{Giaquinto2007}.

\end{document}
